\phantomsection
\label{page:body-second}

\section{研究方法}

为保证旋流分选机理分析的可解释性，本文采用理论分析、实验对比与图表展示相结合的方式组织样例内容。本页保留表格、插图与中英文标题的典型版式，用于后续对齐测试和视觉回归。

\begin{table}[htbp]
    \centering
    \caption{筛分粒度组成\\\textnormal{Table 1-1 Particle size distribution results}}
    \label{tab:sieve}
    \begin{tabular}{C{2.4cm}C{2.0cm}C{2.0cm}C{2.4cm}C{2.4cm}}
        \toprule
        粒级，mm & 产率，\% & 灰分，\% & 累计产率，\% & 累计灰分，\% \\
        \midrule
        $>0.5$ & 3.80 & 7.38 & 3.80 & 7.38 \\
        0.5$\sim$0.25 & 4.55 & 4.56 & 8.35 & 5.84 \\
        0.25$\sim$0.125 & 3.32 & 5.47 & 11.67 & 5.74 \\
        0.125$\sim$0.074 & 4.74 & 3.63 & 16.41 & 5.13 \\
        0.074$\sim$0.045 & 10.72 & 3.11 & 27.13 & 4.33 \\
        $<0.045$ & 72.87 & 4.64 & 100.00 & 4.56 \\
        \bottomrule
    \end{tabular}
\end{table}

\begin{figure}[htbp]
    \centering
    \fbox{\rule{0pt}{5cm}\rule{11cm}{0pt}}
    \caption{循环矿浆压力与柱体背压的关系\\\textnormal{Figure 1-1 Relationship between the pressure of circulating pulp and the back pressure}}
    \label{fig:pressure}
\end{figure}

从 \figref{fig:pressure} 和 \tabref{tab:sieve} 可以看出，参数变化会直接改变流场稳定性与分选指标。后续章节可在这一版式基础上继续扩展完整正文内容。

\clearpage
